% ---------------------------------------------------------------------------
% AJLATT: active joint localisation and target tracking
% ---------------------------------------------------------------------------
\begingroup
% Line breaking for paragraphs with long, unbreakable identifiers (local to
% this chapter): allow looser lines instead of overfull ones.
\setlength{\emergencystretch}{3em}\tolerance=1000
\chapter{AJLATT: Active Joint Localisation and Target Tracking}\label{chap:ajlatt}

In AJLATT a team of mobile robots with range-bearing sensors has to localise
itself and track a moving target on an occupancy grid map. Every robot runs an
estimator for its own pose and for the target; estimates are exchanged within a
limited communication range and fused with covariance intersection. The robots
choose their velocities, and the reward drives the team to keep the target and
its own poses well estimated while avoiding obstacles, the map border and each
other. The environment is used to study active perception and
information-driven control, cooperative localisation and distributed target
tracking, and multi-agent reinforcement learning with continuous actions,
partial observability and safety constraints.

The main entry points are exported by \code{env\_lib.ajlatt\_env}:
\begin{description}[style=nextline,leftmargin=3.9cm]
  \item[\code{AJLATTEnv}] the Gymnasium environment with joint multi-agent
    spaces (also available as \code{env\_lib.AJLATTEnv});
  \item[\code{AJLATTConfig}] a validated dataclass holding every parameter;
  \item[\code{TeamRewardWrapper}] a scalar-reward view for single-agent
    algorithms;
  \item[\code{make}] the backwards-compatible factory accepting the original
    keyword names;
  \item[\code{maps}] occupancy maps, the map builder and map plotting
    (Section~\ref{sec:ajlatt-maps}).
\end{description}

% ---------------------------------------------------------------------------
\section{Scenario and motion model}\label{sec:ajlatt-scenario}

The team consists of \code{num\_robots} robots (default 4) and
\code{num\_targets} targets (default 1). Target 0 moves according to a target
policy (Section~\ref{sec:ajlatt-target}); further targets are static landmarks
whose positions are estimated in 2-D. One step lasts $\Delta t$ =
\code{sampling\_period} (0.5\,s), so the default episode of 120 steps covers
60\,s. Robots and the target follow the unicycle model
\begin{equation}
  \begin{bmatrix} p_x \\ p_y \\ \theta \end{bmatrix}_{k+1}
  = \begin{bmatrix} p_x \\ p_y \\ \theta \end{bmatrix}_{k}
  + \begin{bmatrix} v\,\Delta t \cos\theta_k \\ v\,\Delta t \sin\theta_k \\ \omega\,\Delta t \end{bmatrix},
  \label{eq:ajlatt-unicycle}
\end{equation}
with the heading wrapped to $[-\pi, \pi)$. The true robot poses are propagated
with noisy velocities $v \sim \mathcal{N}(v_{\mathrm{cmd}}, \sigma_v^2)$,
$\omega \sim \mathcal{N}(\omega_{\mathrm{cmd}}, \sigma_\omega^2)$, where
$(\sigma_v, \sigma_\omega)$ = (\code{sigma\_vR}, \code{sigma\_wR}); with
\code{process\_noise\_fixed=False} the noise scales with the speed,
$\sigma_v = \sigma/\sqrt{2}$ and $\sigma_\omega = 2\sqrt{2}\,\sigma$ with
$\sigma$ = \code{sigmapR}$\,\cdot v_{\mathrm{cmd}}$. The target moves exactly
with the command of its policy, while the robots predict their target beliefs
with a noisy copy of that command (\code{sigma\_vT}, \code{sigma\_wT}, or
\code{sigmapT}$\,\cdot$\code{target\_linear\_velocity} in the speed-dependent
case).

% ---------------------------------------------------------------------------
\section{Sensing, communication and estimation}\label{sec:ajlatt-estimation}

\paragraph{Beliefs.} Robot $i$ keeps a Gaussian belief
$(\hat x_i, P_i)$ of its own pose and a belief
$(\hat x^{T}_{i,t}, P^{T}_{i,t})$ of every target $t$ (a pose for target 0, a
position for landmarks). At reset the means are drawn around the initial poses
with the initial covariances: $P_i$ = \code{robot\_init\_cov[i]}$\,I$ with
heading variance $10^{-3}$, and $P^T_{i,t}$ = \code{target\_init\_cov[t]}$\,I$
with heading variance 0.1 for target 0.

\paragraph{Prediction.} Every belief is predicted with an extended Kalman filter
step for \eqref{eq:ajlatt-unicycle}:
\begin{equation}
  \hat x \leftarrow f(\hat x, u), \qquad
  P \leftarrow \Phi P \Phi^\top + G Q G^\top, \qquad
  Q = \operatorname{diag}(\sigma_v^2, \sigma_\omega^2),
\end{equation}
where $\Phi$ is the Jacobian of the motion with respect to the pose and
$G = \Delta t \bigl[\begin{smallmatrix} \cos\theta & 0 \\ \sin\theta & 0 \\ 0 & 1 \end{smallmatrix}\bigr]$
its Jacobian with respect to the velocity noise.

\paragraph{Measurements.} Robot $\ell$ measures the range and bearing
$z = (\rho, \beta)$ of another robot or a target $j$ when
\code{sensor\_r\_min} $< \rho <$ \code{sensor\_r\_max},
$|\beta| \le$ \code{fov}$/2$, the straight line between the two (true)
positions does not cross an occupied cell of the grid (Bresenham ray casting),
and robot $\ell$'s estimated position lies inside the map. The measurement noise
has standard deviations $\sigma_\rho$ = \code{sigma\_p}$\,\cdot\rho$ (with
\code{range\_noise\_proportional=True}) or \code{sigma\_r}, and
$\sigma_\beta$ = \code{sigma\_th}. All three default to zero: by default the
measurements are exact, and the uncertainty of the estimates stems from the
motion noise and from the uncertainty of the measured party's own estimate.

\paragraph{Communication.} Robots closer than \code{commu\_r\_max} (true
distance) are neighbours; the neighbour set $\mathcal{N}_i$ of robot $i$
includes $i$ itself.

\paragraph{Update.} With \code{use\_update=True} (default) the robots are
processed one after another in index order, and each update is applied in place,
so later robots use the updated beliefs of earlier ones:
\begin{enumerate}
  \item \emph{Self-localisation.} Every measurement $z$ of an entity $j$ with
        belief $(\hat x_j, P_j)$ contributes the information pair
        \begin{equation*}
          S = H_i^\top W H_i, \qquad
          y = H_i^\top W \bigl(z - \hat z + H_i \hat x_i\bigr), \qquad
          W = \bigl(R + H_j P_j H_j^\top\bigr)^{-1},
        \end{equation*}
        where $\hat z$ is the predicted measurement and $H_i$, $H_j$ are the
        Jacobians of the range-bearing model with respect to the observer and
        the observed pose. These pairs and the prior
        $(P_i^{-1}, P_i^{-1}\hat x_i)$ are fused by covariance intersection.
  \item \emph{Target tracking.} For each target, the measurements of all
        neighbours $r \in \mathcal{N}_i$ that detect it contribute analogous
        information pairs (with the roles of observer and target exchanged);
        these are fused by covariance intersection in information form. The
        neighbours' prior beliefs of the target are fused in the same way, and
        the two results are combined by a final covariance intersection. If no
        neighbour detects the target, robot $i$'s belief is the fusion of the
        neighbours' priors.
\end{enumerate}
With \code{use\_update=False} the robots only dead-reckon.

\paragraph{Covariance intersection.} Given information pairs
$(S_k, y_k)$, $k = 1, \dots, K$, covariance intersection chooses weights $c$ on
the probability simplex that minimise the trace of the fused covariance,
\begin{equation}
  c^\star = \operatorname*{arg\,min}_{c \ge 0,\; \sum_k c_k = 1}
    \operatorname{tr}\Bigl(\bigl(\textstyle\sum_k c_k S_k\bigr)^{-1}\Bigr),
  \qquad
  P = \Bigl(\sum_k c^\star_k S_k\Bigr)^{-1},\quad
  \hat x = P \sum_k c^\star_k y_k .
\end{equation}
The fused estimate is consistent without knowledge of the correlations between
the sources. The solver is chosen with \code{ci\_solver}:
\code{"newton"} (default) is an active-set projected Newton method with the
analytic gradient $g_k = -\operatorname{tr}(P S_k P)$ and Hessian
$H_{kl} = 2 \operatorname{tr}(P S_k P S_l P)$; its objective is never worse than
that of \code{"slsqp"}, the original solver (SciPy SLSQP with finite-difference
gradients), which reproduces trajectories of the implementation before
version 1.0 to about $10^{-6}$.

% ---------------------------------------------------------------------------
\section{Spaces}\label{sec:ajlatt-spaces}

\paragraph{Action.} \code{Box} of shape \code{(num\_robots, 2)}, one command
$(v, \omega)$ per robot with $0 \le v \le v_{\max}$ and
$|\omega| \le \omega_{\max}$, where $v_{\max}$ =
\code{max\_linear\_velocity} (default 2\,m/s) and $\omega_{\max}$ =
\code{max\_angular\_velocity} (default $\pi/4$\,rad/s). The per-robot space
is \code{env.single\_action\_space}. Actions are not clipped unless
\code{clip\_actions=True} (the original environment did not clip). NumPy arrays
and tensors are accepted; rows beyond \code{num\_robots} are ignored with a
warning, and fewer rows raise \code{ValueError}.

\paragraph{Observation.} \code{Box} of shape \code{(num\_robots, obs\_dim)}
with $\text{obs\_dim} = 7 + 4(n_R - 1) + 5$ for $n_R$ robots (24 for the
default team); the per-robot space is \code{env.single\_observation\_space}.
Row $i$ is computed from robot $i$'s beliefs; positions and headings of other
entities are expressed in robot $i$'s estimated body frame.
\code{env.observation\_layout()} returns the slices of
Table~\ref{tab:ajlatt-obs}.

\begin{table}[htbp]
  \centering\small
  \caption{AJLATT observation of robot $i$ ($n = 6 + 4(n_R - 1)$).}
  \label{tab:ajlatt-obs}
  \begin{tabularx}{\textwidth}{@{}>{\raggedright\arraybackslash}p{0.33\textwidth}l>{\raggedright\arraybackslash}X@{}}
    \toprule
    Key of \code{observation\_layout()} & Columns & Content \\
    \midrule
    \code{target\_relative\_position} & \code{0:2} & believed target position in the robot frame \\
    \code{target\_relative\_heading} & \code{2} & believed target heading minus own heading \\
    \code{target\_velocity} & \code{3:5} & commanded target velocity $(v, \omega)$ of the last step (zero after reset) \\
    \code{target\_cov\_trace} & \code{5} & $\operatorname{tr} P^T_{i,0}$ \\
    \code{neighbours} & \code{6:n} & for every other robot, in index order: relative $x$, relative $y$, relative heading (estimates), trace of its covariance \\
    \code{closest\_obstacle} & \code{n:n+2} & range and bearing of the closest obstacle cell within the field of view and sensor range; $(r_{\max}, \pi)$ if none; zeros while the robot is within \code{obstacle\_sensing\_margin} of the map border \\
    \code{self\_pose} & \code{n+2:n+5} & own pose estimate $(x, y, \theta)$ in the map frame \\
    \code{self\_cov\_trace} & \code{n+5} & $\operatorname{tr} P_i$ \\
    \bottomrule
  \end{tabularx}
\end{table}

% ---------------------------------------------------------------------------
\section{Reward and episode end}\label{sec:ajlatt-reward}

\code{step()} returns the per-robot rewards as an array of shape
\code{(num\_robots,)}. They are built from the tracking cost $c_i$ and the
penalty $p_i$ of robot $i$ in the step,
\begin{equation}
  c_i = w_T \operatorname{tr} P^T_{i,0} + w_R \operatorname{tr} P_i, \qquad
  p_i = \beta_B\, \mathbf{1}_{\mathrm{border}}(i)
        + \beta_O\, \mathbf{1}_{\mathrm{obstacle}}(i)
        + \beta_M\, \mathbf{1}_{\mathrm{mutual}}(i),
  \label{eq:ajlatt-cost}
\end{equation}
and, in the default \code{reward\_mode="cost"} (the original reward),
\begin{equation}
  r_i = -c_i - p_i,
  \label{eq:ajlatt-reward}
\end{equation}
with $w_T$ = \code{target\_cov\_weight} (10), $w_R$ = \code{robot\_cov\_weight}
(5) and the penalties
\begin{itemize}
  \item $\beta_B$ = \code{boundary\_penalty} (20) when the estimated position
        lies within 0.1\,m of the map border or outside the map;
  \item otherwise $\beta_O$ = \code{obstacle\_penalty} (1.3) when an obstacle
        cell is closer to the estimated pose than
        \code{obstacle\_collision\_distance} (0.2\,m, searched in all
        directions up to \code{sensor\_r\_max}); this also sets the robot's
        collision flag;
  \item $\beta_M$ = \code{mutual\_collision\_penalty} (0.5) when another
        robot's estimated position is closer than
        \code{mutual\_collision\_distance} (0.4\,m).
\end{itemize}
The team reward is $\sum_i r_i$ (\code{info["team\_reward"]}), and
\code{info["tracking\_cost"]} holds $c_i$ for evaluations that do not depend
on the reward settings.

\paragraph{Reward settings.} Four configuration fields change the reward
(Table~\ref{tab:ajlatt-reward-settings}); they are applied in the order of the
table.

\begin{table}[htbp]
  \centering\small
  \caption{Reward settings of AJLATT (defaults first; the defaults give the
    original reward \eqref{eq:ajlatt-reward}).}
  \label{tab:ajlatt-reward-settings}
  \begin{tabularx}{\textwidth}{@{}p{4.2cm}>{\raggedright\arraybackslash}X@{}}
    \toprule
    Setting & Effect \\
    \midrule
    \code{reward\_mode="cost"} & $r_i = -c_i - p_i$: every reward is negative. \\
    \code{reward\_mode="bounded"} & $r_i = e^{-c_i/\sigma} - p_i/\sigma$ with
      $\sigma$ = \code{cost\_scale} (20): at most 1, positive while the robot
      tracks well, so a collision that ends the episode forfeits the rewards of
      the remaining steps. \\
    \code{terminate\_on\_collision} (\code{True}) & An obstacle collision ends
      the robot's episode (\code{terminated[i]}); with \code{False} the
      collision is only penalised and the episode continues. \\
    \code{collision\_termination\_penalty} (0) & Subtracted from $r_i$ on the
      step a collision ends robot $i$'s episode. \\
    \code{team\_reward\_weight} $w$ (0) & $r_i \leftarrow (1-w)\, r_i + w\,
      \frac{1}{n_R}\sum_j r_j$: 0 keeps individual rewards, 1 gives every robot
      the team mean; the team sum is unchanged. \\
    \bottomrule
  \end{tabularx}
\end{table}

\paragraph{Choosing the settings.} With the defaults every reward is negative
and the first collision ends the episode (in the team view of
\code{TeamRewardWrapper} and \code{env\_lib.make\_vec}), so a team can raise its
return by colliding early: on 64 seeded episodes random actions reach $-627$
in 16 steps, the classical controller $-7{,}090$ in 115. The registered id
keeps this reward for compatibility, and \code{env-lib describe} prints a
caution. For learning, use \code{reward\_mode="bounded"} (random 39.9,
controller 223.6) or \code{terminate\_on\_collision=False} (random
$-34{,}273$, controller $-7{,}415$), or a
\code{collision\_termination\_penalty} larger than the cost of an episode.

\begin{pycode}
import env_lib

env = env_lib.make("AJLATT-v0", terminate_on_collision=False)   # penalised, not stopped
env = env_lib.make("AJLATT-v0", reward_mode="bounded", team_reward_weight=0.5)
\end{pycode}

\paragraph{Episode end.} \code{terminated} is a boolean array of shape
\code{(num\_robots,)} holding the obstacle collision flags (all \code{False}
with \code{terminate\_on\_collision=False}). If an estimator diverges
(non-finite covariance) all flags are set, a \code{RuntimeWarning} is issued and
\code{info["numerical\_error"]} is \code{True}. \code{truncated} is a
\code{bool}, set after \code{max\_episode\_steps} steps. The environment keeps
simulating after a collision flag; whether a flag ends the episode is left to
the caller (\code{record\_episode} and \code{TeamRewardWrapper} use
\code{any(terminated)}).

% ---------------------------------------------------------------------------
\section{Configuration}\label{sec:ajlatt-config}

\code{AJLATTEnv(config=None, render\_mode=None, **overrides)} takes an
\code{AJLATTConfig} and/or individual fields as keyword arguments, which are
applied on top of \code{config}. Unknown field names raise \code{TypeError}.

\begin{pycode}
from env_lib import AJLATTConfig, AJLATTEnv

env = AJLATTEnv(map_name="obstacles05", num_robots=3)   # fields as keywords
config = AJLATTConfig(map_name="obstacles04", sensor_r_max=5.0, fov=120.0)
env = AJLATTEnv(config, render_mode="rgb_array")
env = AJLATTEnv(config.replace(ci_solver="slsqp"))              # validated copy
\end{pycode}

\code{AJLATTConfig} validates its values on construction (for example
\code{0 <= sensor\_r\_min < sensor\_r\_max}, \code{0 < fov <= 360}, non-negative
noise levels, enough initial poses and covariances for the team) and offers
\code{replace(**changes)}, \code{to\_dict()} and \code{from\_dict(data)}, the
property \code{observation\_dim}, and \code{from\_kwargs(strict=True, **kwargs)},
which also accepts the legacy names of Table~\ref{tab:ajlatt-legacy} (as do
\code{AJLATTEnv}, \code{env\_lib.ajlatt\_env.make} and the callable module
\code{env\_lib.ajlatt\_env}). All
fields are listed in Table~\ref{tab:ajlatt-params}.

\begingroup\small
\setlength{\LTcapwidth}{\textwidth}
\begin{longtable}{@{}>{\raggedright\arraybackslash}p{0.32\textwidth}>{\raggedright\arraybackslash}p{0.18\textwidth}>{\raggedright\arraybackslash}p{0.44\textwidth}@{}}
  \caption{Fields of \code{AJLATTConfig} (keyword arguments of \code{AJLATTEnv}).}\label{tab:ajlatt-params}\\
  \toprule
  Field & Default & Description \\
  \midrule
  \endfirsthead
  \toprule
  Field & Default & Description \\
  \midrule
  \endhead
  \bottomrule
  \endfoot
  \multicolumn{3}{@{}l}{\textit{Scenario}} \\
  \code{map\_name} & \code{"obstacles04"} & Bundled map name or path to a user map (Section~\ref{sec:ajlatt-maps}). \\
  \code{num\_robots} & \code{4} & Team size $n_R$ ($\ge 1$). \\
  \code{num\_targets} & \code{1} & Number of targets; target 0 moves, the others are static landmarks. \\
  \code{max\_episode\_steps} & \code{120} & Episode length; reaching it sets \code{truncated}. \\
  \code{sampling\_period} & \code{0.5} & Duration $\Delta t$ of one step in seconds. \\
  \code{margin2wall} & \code{1.0} & Safety margin of the map queries \code{in\_bound} and \code{is\_collision}. \\
  \code{robot\_init\_pose} & six poses & $(x, y, \theta)$ per robot; the first four defaults are $(7.5, 12.5, 0)$, $(7.5, 10.5, 0)$, $(10, 12.5, 0)$, $(10, 10.5, 0)$. \\
  \code{target\_init\_pose} & 13 poses & $(x, y, \theta)$ per target; target 0 starts at $(10.5, 11.5, 0)$. \\
  \code{robot\_init\_cov} & \code{(0.1,) * 6} & Initial position variance per robot. \\
  \code{target\_init\_cov} & \code{(1.0,) + (10.0,) * 10} & Initial position variance per target. \\
  \multicolumn{3}{@{}l}{\textit{Sensing and communication}} \\
  \code{sensor\_r\_min} & \code{0.0} & Minimum sensing range (m). \\
  \code{sensor\_r\_max} & \code{3.0} & Maximum sensing range (m); also the range of the obstacle reading. \\
  \code{fov} & \code{180.0} & Field of view in degrees, centred on the heading. \\
  \code{sigma\_p} & \code{0.0} & Range noise as a fraction of the range (proportional mode). \\
  \code{sigma\_r} & \code{0.0} & Absolute range noise (m) when \code{range\_noise\_proportional=False}. \\
  \code{sigma\_th} & \code{0.0} & Bearing noise (rad). \\
  \code{range\_noise\_proportional} & \code{True} & Use \code{sigma\_p} instead of \code{sigma\_r}. \\
  \code{use\_update} & \code{True} & Run measurement updates (\code{False}: dead reckoning only). \\
  \code{commu\_r\_max} & \code{6.0} & Communication range (m). \\
  \code{ci\_solver} & \code{"newton"} & Covariance-intersection solver, \code{"newton"} or \code{"slsqp"}. \\
  \multicolumn{3}{@{}l}{\textit{Motion noise}} \\
  \code{process\_noise\_fixed} & \code{True} & Fixed velocity noise; \code{False} scales it with the speed. \\
  \code{sigmapR} & \code{0.1} & Robot noise factor in the speed-dependent mode. \\
  \code{sigma\_vR} & \code{0.1} & Robot linear-velocity noise (m/s). \\
  \code{sigma\_wR} & $0.5\pi/180$ & Robot angular-velocity noise (rad/s, 0.5 degrees/s). \\
  \code{sigmapT} & \code{0.1} & Target noise factor in the speed-dependent mode. \\
  \code{sigma\_vT} & \code{0.2} & Target linear-velocity noise used by the estimators (m/s). \\
  \code{sigma\_wT} & $0.5\pi/180$ & Target angular-velocity noise used by the estimators (rad/s). \\
  \multicolumn{3}{@{}l}{\textit{Target motion}} \\
  \code{target\_policy} & \code{"auto"} & Target policy (Table~\ref{tab:ajlatt-target}). \\
  \code{target\_linear\_velocity} & \code{0.3} & Speed of the \code{"sine"} policy; also scales the speed-dependent target noise. \\
  \code{target\_omega\_bound} & $\pi/4$ & Turn-rate bound of the \code{"random"} policy. \\
  \code{k\_theta} & \code{0.5} & Heading gain of the \code{"sine"} and \code{"waypoints"} policies. \\
  \code{k\_p} & \code{1.0} & Kept for compatibility; currently unused. \\
  \multicolumn{3}{@{}l}{\textit{Actions}} \\
  \code{max\_linear\_velocity} & \code{2.0} & Upper bound of $v$ (m/s); the lower bound is 0. \\
  \code{max\_angular\_velocity} & $\pi/4$ & Bound of $|\omega|$ (rad/s). \\
  \code{clip\_actions} & \code{False} & Clip actions to the action space. \\
  \multicolumn{3}{@{}l}{\textit{Rewards and termination}} \\
  \code{target\_cov\_weight} & \code{10.0} & Weight $w_T$ of the target covariance trace. \\
  \code{robot\_cov\_weight} & \code{5.0} & Weight $w_R$ of the self-localisation covariance trace. \\
  \code{boundary\_penalty} & \code{20.0} & Penalty $\beta_B$ near or beyond the map border (0.1\,m margin). \\
  \code{obstacle\_penalty} & \code{1.3} & Penalty $\beta_O$ for an obstacle closer than \code{obstacle\_collision\_distance}. \\
  \code{obstacle\_collision\_distance} & \code{0.2} & Obstacle collision distance (m). \\
  \code{mutual\_collision\_penalty} & \code{0.5} & Penalty $\beta_M$ for robots closer than \code{mutual\_collision\_distance}. \\
  \code{mutual\_collision\_distance} & \code{0.4} & Robot-robot distance threshold (m). \\
  \code{terminate\_on\_collision} & \code{True} & End a robot's episode (\code{terminated}) at an obstacle collision; \code{False} only penalises it. \\
  \code{reward\_mode} & \code{"cost"} & \code{"cost"} or \code{"bounded"} (Table~\ref{tab:ajlatt-reward-settings}). \\
  \code{cost\_scale} & \code{20.0} & Scale $\sigma$ of the \code{"bounded"} reward. \\
  \code{collision\_termination\_penalty} & \code{0.0} & Extra penalty when a collision ends a robot's episode. \\
  \code{team\_reward\_weight} & \code{0.0} & Weight $w \in [0, 1]$ of the team mean in each robot's reward. \\
  \code{obstacle\_sensing\_margin} & \code{1.0} & Robots closer than this to the map border report a zero obstacle reading (original behaviour). \\
  \multicolumn{3}{@{}l}{\textit{Miscellaneous}} \\
  \code{seed} & \code{None} & Seed of the first \code{reset()} when none is passed. \\
\end{longtable}
\endgroup

The default poses and covariances cover up to six robots and eleven targets;
larger teams need explicit \code{robot\_init\_pose} and \code{robot\_init\_cov}
(and likewise for targets).

\begin{table}[htbp]
  \centering\small
  \caption{Legacy parameter names; each emits a \code{DeprecationWarning}.}
  \label{tab:ajlatt-legacy}
  \begin{tabularx}{\textwidth}{@{}l>{\raggedright\arraybackslash}X@{}}
    \toprule
    Legacy name & Replacement \\
    \midrule
    \code{num\_Robot}, \code{num\_Target} & \code{num\_robots}, \code{num\_targets} \\
    \code{T\_steps} & \code{max\_episode\_steps} \\
    \code{maxmum\_run} & caps \code{max\_episode\_steps} at \code{maxmum\_run + 1} \\
    \code{SIGPERCENT} & \code{range\_noise\_proportional} \\
    \code{useupdate} & \code{use\_update} \\
    \code{fix\_seed} & \code{seed} (\code{fix\_seed=True} sets \code{seed=1}) \\
    \code{render=True} & \code{render\_mode="human"} \\
    obsolete keys (\code{figID}, \code{ssh\_debug}, \code{episode\_num}, \dots) & ignored with a warning \\
    \bottomrule
  \end{tabularx}
\end{table}

\paragraph{Command-line integration.} \code{AJLATTConfig.add\_arguments(parser,
prefix="")} registers every scalar field as an option \code{-{}-<prefix><name>}
on an \code{argparse} parser, and \code{AJLATTConfig.from\_namespace(namespace,
prefix="")} builds a configuration from the parsed arguments. Training scripts
keep their own options next to the environment's:

\begin{pycode}
import argparse
from env_lib import AJLATTConfig, AJLATTEnv

parser = argparse.ArgumentParser()
parser.add_argument("--lr", type=float, default=3e-4)     # training options
AJLATTConfig.add_arguments(parser, prefix="env.")        # --env.map_name, ...
args = parser.parse_args(["--env.map_name", "obstacles05",
                          "--env.num_robots", "3"])
env = AJLATTEnv(AJLATTConfig.from_namespace(args, prefix="env."))
\end{pycode}

% ---------------------------------------------------------------------------
\section{Target motion}\label{sec:ajlatt-target}

The target policy maps the true target pose to a unicycle command. Policies are
per-environment objects that are reset with the environment's generator at
every \code{reset()}.

\begin{table}[htbp]
  \centering\small
  \caption{Target policies (\code{target\_policy}).}
  \label{tab:ajlatt-target}
  \begin{tabularx}{\textwidth}{@{}l>{\raggedright\arraybackslash}X@{}}
    \toprule
    Policy & Behaviour \\
    \midrule
    \code{"auto"} & \code{"sine"} on the map \code{empty}, \code{"waypoints"} on
      \code{obstacles04} and \code{obstacles05}, \code{"circle"} otherwise \\
    \code{"sine"} & follows the slope field $dy/dx = 5\cos x$ at speed
      \code{target\_linear\_velocity} with heading gain \code{k\_theta} until
      $x \ge 32$\,m, then circles with $v = 0.25$\,m/s on a 4\,m radius \\
    \code{"waypoints"} & proportional heading control (gain \code{k\_theta})
      towards piecewise waypoints; built-in routes for \code{obstacles04}
      (0.1\,m/s) and \code{obstacles05} (0.25\,m/s); other routes with
      \code{env\_lib.ajlatt\_env.controllers.WaypointPolicy} \\
    \code{"circle"} & constant $(v, \omega) = (0.25\,\mathrm{m/s}, 0.15\,\mathrm{rad/s})$ \\
    \code{"random"} & $v = 0.25$\,m/s and a turn rate drawn uniformly from
      $[-$\code{target\_omega\_bound}, \code{target\_omega\_bound}$]$ at every step \\
    \code{"static"} & the target does not move \\
    \bottomrule
  \end{tabularx}
\end{table}

% ---------------------------------------------------------------------------
\section{Info dictionary and episode statistics}\label{sec:ajlatt-info}

\begin{table}[htbp]
  \centering\small
  \caption{AJLATT info dictionary ($n_R$ = \code{num\_robots}).}
  \label{tab:ajlatt-info}
  \begin{tabularx}{\textwidth}{@{}l>{\raggedright\arraybackslash}X@{}}
    \toprule
    Key & Content \\
    \midrule
    \code{agent\_rewards} & \code{float64} array $(n_R,)$, equal to the returned rewards (zeros at reset) \\
    \code{team\_reward} & sum of the per-robot rewards \\
    \code{collisions} & boolean array $(n_R,)$: obstacle collisions in this step \\
    \code{episode\_collisions} & integer array $(n_R,)$: number of border, obstacle and robot-robot penalty events in the episode \\
    \code{target\_observed} & boolean array $(n_R,)$: robot measured target 0 in this step \\
    \code{communication} & boolean matrix $(n_R, n_R)$: communication links, including self-loops \\
    \code{robot\_cov\_trace} & array $(n_R,)$: $\operatorname{tr} P_i$ \\
    \code{target\_cov\_trace} & array $(n_R,)$: $\operatorname{tr} P^T_{i,0}$ \\
    \code{tracking\_cost} & array $(n_R,)$: tracking cost $c_i$ of \eqref{eq:ajlatt-cost} (\code{step()} only) \\
    \code{step} & number of steps taken \\
    \code{numerical\_error} & \code{True} if an estimator diverged (\code{step()} only) \\
    \bottomrule
  \end{tabularx}
\end{table}

\code{env.episode\_statistics()} returns per-step traces of the current
episode as arrays of shape \code{(steps + 1, num\_robots)}, starting with the
state after reset: \code{robot\_cov\_trace}, \code{target\_cov\_trace},
\code{robot\_error} and \code{target\_error} (Euclidean position errors of the
robots' own and target beliefs) and \code{reward}. The traces are cleared at
every reset.

% ---------------------------------------------------------------------------
\section{Team reward wrapper and legacy factory}\label{sec:ajlatt-wrapper}

\code{TeamRewardWrapper(env)} returns the team reward $\sum_i r_i$ as a
\code{float} and \code{any(terminated)} as a \code{bool}, for algorithms that
expect scalar values. The per-robot values remain available in the info
dictionary under \code{"agent\_rewards"} and \code{"agent\_terminated"}.

\begin{pycode}
from env_lib import AJLATTEnv
from env_lib.ajlatt_env import TeamRewardWrapper

env = TeamRewardWrapper(AJLATTEnv(map_name="obstacles04"))
obs, info = env.reset(seed=1)
obs, reward, terminated, truncated, info = env.step(env.action_space.sample())
\end{pycode}

The original factory remains available for existing code:
\code{env\_lib.ajlatt\_env.make(**kwargs)}, or the callable module
\code{env\_lib.ajlatt\_env(num\_Robot=4, T\_steps=100)}. It translates the legacy
names of Table~\ref{tab:ajlatt-legacy}, ignores unknown keys with a warning and
maps \code{render=True} to \code{render\_mode="human"}. New code should create
\code{AJLATTEnv} directly.

% ---------------------------------------------------------------------------
\section{Maps}\label{sec:ajlatt-maps}

The subpackage \code{env\_lib.ajlatt\_env.maps} contains the occupancy maps and
the tools to create and inspect them.

\subsection{Bundled maps}\label{sec:ajlatt-bundled-maps}

\code{available\_maps()} lists the bundled maps (Table~\ref{tab:ajlatt-maps});
\code{ajlatt-plot-map -{}-list} prints the same list. All maps start at the
origin (\code{mapmin} $= (0, 0)$). Figure~\ref{fig:ajlatt-map} shows
\code{obstacles05}.

\begin{table}[htbp]
  \centering\small
  \caption{Bundled maps. Size in cells ($n_x \times n_y$) and metres.}
  \label{tab:ajlatt-maps}
  \begin{tabularx}{\textwidth}{@{}>{\raggedright\arraybackslash}p{0.34\textwidth}lll>{\raggedright\arraybackslash}X@{}}
    \toprule
    Map & Cells & Cell size & Extent & Notes \\
    \midrule
    \code{obstacles} & $51 \times 51$ & 0.5\,m & 25.5\,m & \\
    \code{obstacles02}, \code{obstacles07}, \code{obstacles07-02}, \code{obstacles08} & $181 \times 181$ & 0.2\,m & 36.2\,m & \\
    \code{obstacles03} & $102 \times 102$ & 0.2\,m & 20.4\,m & \\
    \code{obstacles04} & $181 \times 181$ & 0.2\,m & 36.2\,m & default; waypoint route \\
    \code{obstacles05} & $181 \times 181$ & 0.4\,m & 72.4\,m & waypoint route \\
    \code{obstacles06}, \code{obstacles06-02}, \code{obstacles06-03}, \code{obstacles08-02} & $181 \times 181$ & 0.1\,m & 18.1\,m & \\
    \code{example\_complete\_map}, \code{sample\_map\_config} & $181 \times 181$ & 0.4\,m & 72.4\,m & map builder examples \\
    \code{empty} & $82 \times 62$ & 0.5\,m & $41 \times 31$\,m & no obstacles; sine target policy \\
    \code{empty05} & $181 \times 181$ & 0.4\,m & 72.4\,m & no obstacles \\
    \code{emptySmall} & $135 \times 135$ & 0.2\,m & 27\,m & no obstacles \\
    \code{emptySmall05} & $65 \times 65$ & 0.4\,m & 26\,m & no obstacles \\
    \code{dynamic\_map}, \code{dynamic\_map\_2} & $181 \times 181$ & 0.4\,m & 72.4\,m & obstacles resampled at every reset \\
    \bottomrule
  \end{tabularx}
\end{table}

\begin{figure}[htbp]
  \centering
  \includegraphics[width=0.55\textwidth]{map_obstacles05.png}
  \caption{The bundled map \code{obstacles05} ($72.4 \times 72.4$\,m,
    0.4\,m cells) plotted with \code{ajlatt-plot-map obstacles05}; occupied
    cells are grey, including the one-cell wall around the map.}
  \label{fig:ajlatt-map}
\end{figure}

\subsection{Map files}\label{sec:ajlatt-map-format}

A static map consists of two files with the same name:
\begin{itemize}
  \item \code{<name>.yaml}, the header, with the keys \code{mapdim}
        ($[n_x, n_y]$, cells along $x$ and $y$), \code{mapres} (cell size
        $[r_x, r_y]$ in metres), \code{mapmin} and \code{mapmax} (corners of the
        map in metres) and optionally \code{origin}. Further keys written by
        older tools or the map builder (\code{datatype}, \code{mappath},
        \code{origincells}, \code{storage}, the embedded specification) are
        ignored by the loader.
  \item \code{<name>.cfg}, the occupancy grid as whitespace-separated numbers,
        one text row per grid row: $n_y$ rows of $n_x$ values, where row 0 is
        the bottom row ($y$ = \code{mapmin}) and 1 marks an occupied cell.
\end{itemize}
Cell $(c_x, c_y)$ has its centre at $(c + 0.5)\,r + $\code{mapmin}. A grid with
fewer cells than its header declares is padded with occupied cells and a
warning. An empty (zero-byte) \code{.cfg} file denotes an obstacle-free map in
which only the boundary blocks rays; a map whose name contains \code{empty} may
omit the file altogether, while a missing grid for any other static map raises
\code{FileNotFoundError}. A dynamic map has only a header, which defines
\code{submaporigin} (four $(\text{row}, \text{column})$ cell anchors, flattened
to eight numbers) and \code{lib\_path} (a directory of \code{.npy} obstacle
masks relative to the header).

To use your own map, pass its path to \code{map\_name}, with or without the
\code{.yaml} suffix, for example \code{AJLATTEnv(map\_name="maps/my\_map")};
the grid file must sit next to the header. Check that the default initial poses
(Table~\ref{tab:ajlatt-params}) are free cells of the map, or pass
\code{robot\_init\_pose} and \code{target\_init\_pose}.

\subsection{Map API}\label{sec:ajlatt-map-api}

\begin{itemize}
  \item \code{load\_grid\_map(name\_or\_path, margin2wall=0.5)} returns a
        \code{GridMap}, or a \code{DynamicMap} for a header with
        \code{submaporigin} and no grid. \code{load\_map(name)} returns
        \code{(occupancy, mapmin, mapmax)} for plotting, and
        \code{resolve\_map\_path(name)} the file path without suffix.
  \item \code{GridMap} holds the grid as \code{map} (shape $(n_y, n_x)$,
        \code{None} for empty maps) together with \code{mapdim}, \code{mapres},
        \code{mapmin} and \code{mapmax}; its queries are listed in
        Table~\ref{tab:ajlatt-gridmap}. \code{GridMap.from\_array(occupancy,
        mapmin=(0, 0), mapres=0.4)} builds a map from an array of shape
        $(n_y, n_x)$.
  \item \code{DynamicMap.generate\_map(chosen\_idx=None, rot\_angs=None,
        rng=None)} places four distinct obstacles from its library at the four
        anchors, each rotated by a random multiple of 18 degrees, and closes the
        border. The environment calls it with its own generator at every reset.
\end{itemize}

\begin{table}[htbp]
  \centering\small
  \caption{Queries of \code{GridMap} (positions and poses in metres and radians).}
  \label{tab:ajlatt-gridmap}
  \begin{tabularx}{\textwidth}{@{}>{\raggedright\arraybackslash}p{0.47\textwidth}>{\raggedright\arraybackslash}X@{}}
    \toprule
    Method & Result \\
    \midrule
    \code{is\_blocked(start, end)} & \code{True} if the straight line crosses an occupied
      cell (exact Bresenham ray casting) \\
    \code{is\_blocked\_batch(starts, ends)} & the same for arrays of shape \code{(k, 2)} or \code{(k, 3)} \\
    \code{get\_closest\_obstacle(pose, ang\_res=0.05, fov=pi, r\_max=3.0)} & range and
      bearing of the closest occupied cell in the field of view, or \code{None} \\
    \code{is\_collision(pos, margin=None)} & \code{True} if \code{pos} is out of bounds or
      within \code{margin} (default \code{margin2wall}) of an occupied cell \\
    \code{in\_bound(pos)} & \code{True} if \code{pos} is at least \code{margin2wall} inside the map \\
    \code{se2\_to\_cell(pos)}, \code{cell\_to\_se2(cell)} & conversion between positions and cell indices \\
    \bottomrule
  \end{tabularx}
\end{table}

\subsection{Building maps}\label{sec:ajlatt-builder}

The map builder (\code{env\_lib.ajlatt\_env.maps.builder}, command
\code{ajlatt-build-map}) rasterises a YAML specification that describes the map
in metres:

\begin{textcode}
map_info:
  name: my_map
  width: 36.0            # metres along x
  height: 36.0           # metres along y
  resolution: 0.2        # cell size in metres
  boundary:
    enabled: true
    thickness: 0.2
    style: rectangle     # rectangle | circle | polygon
    rectangle: {margin: 0.0}
obstacles:
  - {type: rectangle, center: [10, 10], width: 4, height: 2, angle: 0.3}
  - {type: circle, center: [25, 25], radius: 3, filled: false}
  - {type: polygon, points: [[5, 20], [9, 20], [7, 26]]}
  - {type: line, start: [0, 18], end: [12, 18], width: 0.4}
  - {type: random, num_obstacles: 5, min_size: 0.5, max_size: 1.5, seed: 3}
\end{textcode}

The grid has $n_x = \operatorname{round}(\text{width}/\text{resolution})$ columns
and $n_y = \operatorname{round}(\text{height}/\text{resolution})$ rows; a cell
is occupied when its centre lies inside a shape. The boundary (enabled by
default) is a rectangular wall with an optional \code{margin}, the outside of a
circle (\code{circle: \{center, radius\}}) or the outside of a polygon
(\code{polygon: \{points\}}), \code{thickness} metres thick. The obstacle types
are listed in Table~\ref{tab:ajlatt-obstacles}; closed shapes accept
\code{filled: false} for a one-cell outline.

\begin{table}[htbp]
  \centering\small
  \caption{Obstacle types of the map builder (lengths in metres, angles in radians).}
  \label{tab:ajlatt-obstacles}
  \begin{tabularx}{\textwidth}{@{}l>{\raggedright\arraybackslash}X@{}}
    \toprule
    Type & Keys \\
    \midrule
    \code{rectangle} & \code{center}, \code{width}, \code{height}, \code{angle} (default 0), \code{filled} \\
    \code{circle} & \code{center}, \code{radius}, \code{filled} \\
    \code{triangle} & \code{vertices} (three points), \code{filled} \\
    \code{hexagon} & \code{center}, \code{radius}, \code{angle}, \code{filled} \\
    \code{polygon} & \code{points} (at least three), \code{filled} \\
    \code{line} & \code{start}, \code{end}, \code{width} (default: one cell) \\
    \code{maze} & \code{cell\_size} (5.0), \code{wall\_thickness} (0.5) and \code{gaps}, a list of
      \code{\{type: horizontal|vertical, position: <cell index>\}} walls to leave open \\
    \code{random} & \code{num\_obstacles} (10), \code{min\_size} (1.0), \code{max\_size} (3.0), \code{seed}: axis-aligned squares \\
    \bottomrule
  \end{tabularx}
\end{table}

The builder writes \code{<name>.yaml} (the grid header followed by the
specification, so a map can always be rebuilt from its own header) and
\code{<name>.cfg}, in the orientation the environment reads. The name is
\code{map\_info.name} unless \code{-n} is given:

\begin{shellcode}
ajlatt-build-map --create-sample my_map        # writes my_map.spec.yaml
ajlatt-build-map my_map.spec.yaml --png        # my_map.yaml, .cfg, .png
ajlatt-build-map spec.yaml -o maps/ -n arena   # other directory and name
\end{shellcode}

\code{-{}-create-sample} sets \code{map\_info.name} to the requested file name,
so the second command writes \code{my\_map.yaml}, \code{my\_map.cfg} and
\code{my\_map.png}. The same functionality is available from Python:
\code{build\_map(spec, output\_dir=None, *, name=None, png=False)} accepts a path
or a dictionary and returns the paths of the header and grid,
\code{rasterize(spec)} returns the occupancy array and the header, and
\code{save\_map} writes them.

\begin{pycode}
from env_lib import AJLATTEnv
from env_lib.ajlatt_env.maps.builder import build_map

spec = {
    "map_info": {"name": "corridor", "width": 20.0, "height": 10.0,
                 "resolution": 0.2},
    "obstacles": [
        {"type": "line", "start": [10, 0], "end": [10, 6], "width": 0.4},
        {"type": "circle", "center": [15, 5], "radius": 1.5},
    ],
}
header, grid = build_map(spec, "maps", png=True)   # maps/corridor.yaml, ...
env = AJLATTEnv(map_name="maps/corridor",
                robot_init_pose=[(2, 2, 0), (2, 4, 0), (4, 2, 0), (4, 4, 0)],
                target_init_pose=[(5, 5, 0)], target_policy="static")
\end{pycode}

\subsection{Plotting maps}\label{sec:ajlatt-plot}

\code{ajlatt-plot-map} draws a bundled or user map; without \code{-{}-save} it
writes \code{<map>.png} in the working directory. It uses the light theme unless
\code{-{}-theme dark} is given. Dynamic maps are drawn with a fixed obstacle
sample.

\begin{shellcode}
ajlatt-plot-map --list                          # bundled map names
ajlatt-plot-map obstacles04 --save o4.png
ajlatt-plot-map maps/corridor --theme dark
\end{shellcode}

From Python, \code{plot\_map(map\_or\_name, ax=None, *, theme=None,
title=None, save=None, dpi=150)} in \code{env\_lib.ajlatt\_env.maps.plotting}
draws a map into existing axes (or a new off-screen figure) and returns the
axes.

% ---------------------------------------------------------------------------
\section{Rendering}\label{sec:ajlatt-render}

The dashboard ($1100 \times 620$ pixels, Figure~\ref{fig:ajlatt}) shows the map
on the left with the robots' true poses as oriented arrowheads, their pose
beliefs as covariance ellipses, their sensing sectors, fading trajectory trails,
communication links, the active target measurements, the true target (star)
and every robot's belief about the target (crosses). The two panels on the
right plot the trace of each robot's target covariance and of its
self-localisation covariance over the episode. The header shows the map, the
step, the simulated time, the team reward and how many robots currently see
the target. Human mode requires an interactive matplotlib backend; TkAgg needs
the \pkg{tk} bindings of the Python installation (\code{python3-tk} on Debian
and Ubuntu, \code{tk} in conda).

\begin{figure}[htbp]
  \centering
  \includegraphics[width=\textwidth]{ajlatt.png}
  \caption{AJLATT dashboard on \code{obstacles04} (four robots running the
    encircling heuristic of the example script, step 90 of 120). The robots
    surround the target; the right panels show the trace of the target
    covariance and of the self-localisation covariances of the four robots.}
  \label{fig:ajlatt}
\end{figure}

% ---------------------------------------------------------------------------
\section{Registered id}\label{sec:ajlatt-ids}

\envid{AJLATT-v0} creates \code{AJLATTEnv} with the default configuration (four
robots on \code{obstacles04}). Gymnasium's passive environment checker is
disabled for this id because rewards and \code{terminated} are per-robot
arrays. Keyword arguments of \code{env\_lib.make} are configuration fields, for
example \code{env\_lib.make("AJLATT-v0", map\_name="obstacles05")}.

% ---------------------------------------------------------------------------
\section{Usage}\label{sec:ajlatt-usage}

The following controller sends robot $i$ to the $i$-th of $n_R$ evenly spaced
slots on a circle of radius 1.8\,m around its belief of the target, using only
the observation:

\begin{pycode}
import numpy as np
from env_lib import AJLATTEnv


def encircle(obs, layout, radius=1.8, gain=1.2):
    """Send robot i to the i-th slot of a circle around its target belief."""
    n = len(obs)
    pose = obs[:, layout["self_pose"]].astype(np.float64)   # x, y, theta
    rel = obs[:, layout["target_relative_position"]].astype(np.float64)
    c, s = np.cos(pose[:, 2]), np.sin(pose[:, 2])
    target = pose[:, :2] + np.column_stack([c * rel[:, 0] - s * rel[:, 1],
                                            s * rel[:, 0] + c * rel[:, 1]])
    angle = 2 * np.pi * np.arange(n) / n
    slot = target + radius * np.column_stack([np.cos(angle), np.sin(angle)])
    delta = slot - pose[:, :2]
    heading = np.arctan2(delta[:, 1], delta[:, 0])
    error = np.angle(np.exp(1j * (heading - pose[:, 2])))  # wrapped
    speed = np.minimum(2.0, 0.5 * np.linalg.norm(delta, axis=1))
    v = speed * np.maximum(0.0, np.cos(error))
    w = np.clip(gain * error, -np.pi / 4, np.pi / 4)
    return np.column_stack([v, w])


env = AJLATTEnv(map_name="obstacles04", num_robots=4)
layout = env.observation_layout()
obs, info = env.reset(seed=0)                  # obs: (4, 24) float32
episode_return = np.zeros(env.num_robots)
while True:
    obs, reward, terminated, truncated, info = env.step(encircle(obs, layout))
    episode_return += reward                   # per-robot rewards
    if truncated:                              # collision flags ignored
        break
stats = env.episode_statistics()               # arrays of shape (121, 4)
print(episode_return.sum(), info["episode_collisions"])
print(stats["target_cov_trace"][-1])
env.close()
\end{pycode}

The example script compares a similar heuristic with random actions and can
record an episode:

\begin{shellcode}
python examples/environments/ajlatt_demo.py --episodes 3
python examples/environments/ajlatt_demo.py --save renders/ajlatt.gif
python examples/environments/ajlatt_demo.py --map obstacles05 --render-mode human
\end{shellcode}

With the default configuration the heuristic of the script obtains a team
return of about $-7250$ per episode without collision events and keeps the
trace of the target covariance near 2, whereas random actions obtain returns of
$-25\,000$ or lower with dozens of collision events and target covariance traces
of 7 to 12 (the random actions are not seeded, so these values vary from run to
run).

% ---------------------------------------------------------------------------
\section{Reproducibility and performance}\label{sec:ajlatt-perf}

All randomness (initial beliefs, motion and measurement noise, the random target
policy and dynamic maps) is drawn from \code{self.np\_random}. If
\code{reset()} is called without a seed for the first time and
\code{config.seed} is set, that seed is used. With
\code{ci\_solver="slsqp"} trajectories match the implementation before
version 1.0 to about $10^{-6}$; the default Newton solver finds weights that
are at least as good and therefore produces slightly different, not worse,
estimates.

Measured on the development machine without rendering, a step takes 5--9\,ms on
\code{obstacles04} (90\,ms before 1.0) and 6--7\,ms on \code{obstacles05}
(67\,ms before). The gains come from exact batched Bresenham ray casting
(closest-obstacle queries are 7 to 19 times faster) and from the Newton
covariance-intersection solver. Rendering a frame takes about 25\,ms (about
110\,ms before). The cost of a step grows with the number of robots, because
every robot fuses the measurements and priors of its neighbours.

% ---------------------------------------------------------------------------
\section{Changes in 1.0}\label{sec:ajlatt-changes}

AJLATT was rewritten as a Gymnasium environment. See Appendix~\ref{app:migration}
for migration advice.

\begin{itemize}
  \item The environment class \code{AJLATTEnv} with its configuration
        dataclass \code{AJLATTConfig} replaces the former factory, which parsed
        the command line internally and therefore failed in scripts with their
        own command-line arguments. The legacy factory
        still works and translates the old names with a
        \code{DeprecationWarning}. Training options that used to live in the
        parameter module (PPO, LLM and debug settings) were removed; use
        \code{AJLATTConfig.add\_arguments}.
  \item \code{reset()} returns \code{(observation, info)}. Rendering is off by
        default; pass \code{render\_mode}.
  \item The spaces are joint multi-agent spaces with one row per robot, and
        rewards and \code{terminated} are per-robot arrays; the per-robot
        spaces are available as \code{single\_action\_space} and
        \code{single\_observation\_space}.
  \item The target-velocity entries of the observation contain the commanded
        target velocity instead of a configured constant.
  \item Fixed a memory leak: covariance traces accumulated across episodes
        (per-episode traces are now returned by \code{episode\_statistics()}).
        Also fixed: target policies were shared between environments and never
        reset, unknown maps triggered an interactive prompt, and dynamic maps
        did not work; they now resample at every reset and no longer need
        scikit-image. Three bundled maps that were one row short of their header
        were padded with a wall row.
  \item About ten times faster (Section~\ref{sec:ajlatt-perf}) and a new
        renderer.
  \item Map tools were consolidated into \code{env\_lib.ajlatt\_env.maps}; the
        builder now writes grids in the orientation the environment reads (the
        old builder produced transposed maps) and embeds its specification in
        the header.
  \item The id \envid{AJLATT-v1}, an unusable alias of \envid{AJLATT-v0}, was
        removed.
\end{itemize}

\endgroup
